\chapter*{Agradecimientos}
% [Por completar por los autores.]

\chapter*{Resumen}
\primerparrafo Se estudia numéricamente la inestabilidad de Kelvin--Helmholtz (KHI) en plasmas relativistas dentro del marco de la magnetohidrodinámica relativista resistiva (RRMHD), con el objetivo de caracterizar cómo la resistividad finita modifica el desarrollo de la inestabilidad. Las simulaciones se realizan con el código \textit{Cueva}, que integra las ecuaciones de la RRMHD mediante un esquema de volúmenes finitos de alta resolución (reconstrucción de quinto orden MP5, solucionador de Riemann HLL y limpieza de divergencias GLM), con integración temporal Implícito--Explícita (IMEX-RK) para tratar la rigidez del término resistivo. Sobre un barrido de $44$ simulaciones en la conductividad eléctrica $\sigma\in[100,10500]$ se extrae la tasa de crecimiento lineal $\gKHI(\sigma)$, cuya transición resistivo$\to$ideal se describe con un modelo sigmoide con \emph{offset}; su asíntota ideal, $C+\gamma_0=1.03\pm0.05$, resulta consistente dentro del $8\%$ con la predicción de la teoría lineal RMHD compresible. Se estima una conductividad crítica $\scrit\approx2861$ con una incertidumbre sistemática $\pm1944$. El análisis de las variables globales y no lineales (corriente máxima, disipación, energía magnética, enstrofía y fracción fuera del plano) revela que la conductividad controla la intensidad de las láminas de corriente y el desarrollo de estructuras secundarias, con una firma de reconexión magnética ($\Omaxzp\propto\sigma^{1.22}$) compatible con el escalamiento de Sweet--Parker y el modo \emph{tearing} \sebnota{no olvidar implementar la campaña de campo magnetico}.\claudenota{REVISIÓN GLOBAL (Claude): (1) el Resumen NO menciona la campaña de orientación/intensidad de B (B/C), que sí está en el cap.~6 $\Rightarrow$ añadir una frase. (2) ``Ω$\propto\sigma^{1.22}$ compatible con Sweet--Parker y tearing'' se enuncia como hecho; el exponente es robusto pero el mecanismo (tearing) es \emph{hipótesis} no demostrada $\Rightarrow$ suavizar. (3) liderar con la banda $\scrit\in[2815,6800]$, no con ``$\approx2861$''.}

\noindent\textbf{Palabras clave:} Kelvin--Helmholtz; RRMHD; resistividad; reconexión magnética; tasa de crecimiento; métodos numéricos.

% [Opcional: versión en inglés (Abstract) según formato UD.]

\chapter*{Lista de Figuras}
\listoffigures

\chapter*{Lista de Tablas}
\listoftables

\chapter*{Glosario de Símbolos y Acrónimos}
\begin{description}[leftmargin=3.2cm,style=nextline]
\item[RRMHD] Magnetohidrodinámica Relativista Resistiva.
\item[RMHD] Magnetohidrodinámica Relativista (ideal).
\item[KHI] Inestabilidad de Kelvin--Helmholtz.
\item[GLM] Multiplicadores de Lagrange Generalizados (limpieza de divergencias).
\item[EoS] Ecuación de Estado.
\item[IMEX-RK] Runge--Kutta Implícito--Explícito.
\item[HLL] Solucionador de Riemann Harten--Lax--van Leer.
\item[MP5] Reconstrucción \textit{Monotonicity-Preserving} de quinto orden.
\item[$\gKHI$] Tasa de crecimiento lineal de la KHI.
\item[$\scrit$] Conductividad crítica de transición resistivo--ideal.
\item[$\Rmst$] Número de Reynolds magnético ($\vsh\,\akh\,\sigma$).
\item[$\Ozp$] Enstrofía de perturbación.
\end{description}
% [Glosario por completar/ajustar al formato exigido.]
